\begin{lemma}\label{lem:diamond-descent}
Let $x$ be a non-CM closed point of degree $d$ on $X_1(N)$.
Suppose that $\langle a\rangle$ preserves $x$ and has projective order $r$.
For $M\mid N$, let $s$ be its projective order modulo $M$.
The image of $x$ on $X_1(M)$ is defined over a field of degree dividing
$ds/r$.
\end{lemma}
\begin{proof}
The diamond induces an automorphism of the residue field of $x$.
Its order is $r$, since the action on non-CM points is free.
The subgroup generated by its $s$th power has order $r/s$ and fixes
the image on $X_1(M)$. That image is defined over its fixed field,
whose degree is $ds/r$.
\end{proof}

\begin{proposition}\label{prop:descent31}
Neither $341$ nor $403$ is an order of a torsion point over a decic field.
\end{proposition}
\begin{proof}
At $N=341$ and $403$, the archived complete winding computations
show that $(\langle a\rangle-1)J_1(N)(\mathbb Q)$ is finite for
$a=32$ and $92$, respectively. The character-kernel check is replayed
in \texttt{descent31.m}. In both cases $k_3=8$, and
\[
 80<\frac{325I_N}{32768}.
\]
The minimum CM degrees are $300$ and $120$, respectively.
Thus the global Hecke argument preserves the trace divisor of a
hypothetical decic point under this diamond.

The quintic classification excludes $31$-torsion. Hence its
$31$-level image, and therefore the original moduli point, has exact
degree $10$: a proper residue-field degree would divide $10$ and
would already contradict the known classifications in degrees $1,2,5$.
Both diamonds have order $2$ and reduce to $\pm1$ modulo $31$.
Lemma~\ref{lem:diamond-descent} puts the $31$-level image over a
quintic field, again contradicting the quintic classification.
\end{proof}
